\setlength{\epigraphwidth}{0.6\textwidth}
\chapter{Interlúdio sobre crítica ontológica}

\epigraph{A crítica da economia política é também uma crítica do capitalismo e, para que essa crítica não seja aprisionada no espaço de interpretação que é delimitado pela própria imagem que a sociedade capitalista projeta de si, ela não poderia senão partir do exame crítico das próprias categorias. \parencite[p. 109]{Medeiros2021a}}

\hypertarget{o-significado-de-cruxedtica-ontoluxf3gica}{%
\section{O significado de ``crítica ontológica''}\label{o-significado-de-cruxedtica-ontoluxf3gica}}

Como o presente trabalho se propõe a fazer uma crítica ontológica do assim chamado ``postulado de Khazzoom-Brookes'', nada mais justo do que explicitar o significado de ``crítica'' aqui empregado. \emph{Vide} \textcite{Medeiros2013}, há três significados simultâneos deste termo:\footnote{Também denominado, dentro da tradição filosófica do Realismo Crítico, como ``crítica explanatória'', cf. \textcite{Lawson1997}.}

\begin{quote}
\begin{enumerate}
\item
  a demonstração da falsidade das crenças ou teorias criticadas;
\item
  a simultânea apresentação de uma explicação alternativa e mais abrangente da causalidade de fenômenos anteriormente significados por meio das crenças ou teorias em questão;
\item
  a indicação dos motivos reais que levam à produção e sustentação das concepções equivocadas, mistificadas e/ou ilusórias e, ainda, das condições sociais que facultam a própria crítica. \parencite[p. 78]{Medeiros2013}
\end{enumerate}
\end{quote}

O objetivo da crítica ontológica é elucidar os porquês dos fenômenos analisados, ``atravessar a superfície'' e penetrar ``a variedade do garrido tumulto dos acontecimentos''\footnote{\textcite{Hegel1995}, p.~34. \hl{\textit{Note to self}: Provavelmente vou tirar essa citação daqui e deixar só no final do capítulo.}}, e elucidar a \emph{essência} de cuja \emph{aparência necessária} estes fenômenos são.

Isso permite compreender, em particular, o mérito de \textcite{Khazzoom1980,Khazzoom1987} em elucidar como o próprio paradigma neoclássico já era capaz de explicar o efeito \emph{rebound} dentro de seu próprio arcabouço teórico, não necessitando de argumentos \emph{ad hoc} para contemplar teoricamente esta faceta da realidade empírica. No entanto, ao restringir-se a este paradigma, ele não poderia senão restringir-se aos axiomas neoclássicos de agentes racionais e otimizadores, e, dessa forma, não poderia contemplar as mais fundamentais contradições imbricadas em seu objeto de análise, justamente por abdicar da explicação de seu \emph{explanans}. Em suma, necessariamente teve de assumir a ontologia pressuposta pela economia neoclássica: uma ``ontologia empirista'' na qual somente os fenômenos empíricos são ``reais''.\footnote{Sobre a tendência da Economia a plasmar a realidade social à esfera empírica, Medeiros diz: ``Em lugar de assumir a inspeção crítica das categorias empregadas na vida cotidiana (ou em teorias prévias) como ponto de partida do processo de conhecimento, o pensamento {[}econômico{]} vulgar afunda-se mais e mais na tentativa de aplicar técnicas crescentemente sofisticadas para polir e sistematizar as categorias dadas na prática imediata. Assim segue o pensamento vulgar, com sua atitude indelevelmente passiva e acrítica, a intensa dinâmica do dia a dia. Como consequência, não consegue evitar a eterna ansiedade empirista diante da infinitude da experiência humana, que, no final das contas, bloqueia a busca por estruturas causadoras de fenômenos. Tampouco consegue se livrar de mistificações reais originadas pela sustentação de categorias falsas ou ilusórias em si mesmas.'' \parencite[p. 66]{Medeiros2013}. Ademais: ``Como {[}\ldots{]} a sociedade é um objeto caracterizado pela dinâmica, pela processualidade, pela história, a naturalização de uma determinada forma social exige ações e concepções destinadas a `congelar' a dinâmica social, afastando elementos que possam modificar de forma substantiva este ordenamento particular da sociedade. Por essa razão, pode-se afirmar que a naturalização da forma social corrente é não apenas uma tese ontológica de corte radical, mas também o desdobramento prático da visão de mundo empirista.'' \parencite[p. 102]{Medeiros2013}}

A busca intransigente pela causa genuína do fenômeno presentemente abordado, i.e.~o postulado de Khazzoom-Brookes, requer, necessariamente, uma apreensão aderente do \textit{modo de produção capitalista}, justamente porque aquele fenômeno é por este engendrado --- é dele uma consequência (a se demonstrar) \emph{necessária}. Renegar este elo causal é acanhar-se de compreender a essência verdadeira desse fenômeno; é naturalizar seus pressupostos e, portanto, obscurecer aquilo que se buscava esclarecer. A crítica ontológica, ao contrário, busca aferir tais pressupostos e, ao criticá-los, desnaturalizá-los.

Deve-se tomar em conta, portanto, a totalidade social na qual o capital é ``sujeito automático'', que se torna capaz de pôr seus próprios pressupostos e, neste agir, acabar por moldar o mundo dos homens à sua imagem. Dessa forma, a compreensão de todo fenômeno socioeconômico que diga respeito ao comportamento do capital deve ser visto como ``parte de um complexo dinâmico em interação com outros complexos, como algo que é determinado, tanto interna quanto externamente, por múltiplas leis'' \parencite[p. 338]{Lukacs2012}, justamente porquanto é fenômeno decorrente desta totalidade social. É por esta via que o postulado de Khazzoom-Brookes será abordado neste trabalho.

\hypertarget{o-papel-da-abstrauxe7uxe3o}{%
\section{O papel da abstração}\label{o-papel-da-abstrauxe7uxe3o}}

No prefácio à primeira edição do primeiro volume d'\emph{O Capital}, Marx anuncia o objetivo de sua obra: investigar ``o modo de produção capitalista e suas correspondentes relações de produção e de circulação'' \parencite[p. 78]{Marx2017}. Por não poder, como os cientistas naturais, criar condições experimentais ``que assegur{[}e{]}m o transcurso puro do processo'' a fim de averiguar leis naturais específicas, o cientista social tem de, sempre que possível, aferir tais processos onde e quando eles ``aparecem mais nitidamente e menos obscurecidos por influências perturbadoras''.

Contudo, e nisso Marx destaca-se efetivamente da economia política clássica, é indevido restringir-se à esfera empírica na qual tais fenômenos apresentam-se. É por isso que Marx aborda a ``força da abstração {[}\emph{Abstraktionskraft}{]}'' como o substituto do ``microscópio'' e dos ``reagentes químicos'' dos cientistas sociais: a abstração é o procedimento necessário para abordar (em particular) a crítica do modo de produção capitalista. Isso é esclarecido pelo próprio Marx no prefácio da segunda edição dessa mesma obra:

\begin{quote}
A investigação tem de se apropriar da matéria {[}\emph{Stoff}{]} em seus detalhes, analisar suas diferentes formas de desenvolvimento e \emph{rastrear seu nexo interno}. Somente depois de consumado tal trabalho é que se pode expor adequadamente o movimento real. \parencite[p. 90, grifo meu]{Marx2017}
\end{quote}

O objeto de Marx, o ``modo de produção capitalista'', possui caráter de totalidade social, ``complexo de complexos'', o que requer que sua apreensão científica seja guiada pelo caráter imanente deste objeto. A compreensão de um complexo específico dentro da totalidade social --- por exemplo a esfera da produção de capital, objeto do Livro I d'\emph{O Capital} --- requer que se ``desfoque'' a análise do resto da totalidade à qual este complexo pertence, em prol de sua análise\footnote{No sentido etimológico original do termo: o sufixo ``lise'' remete à ideia de ``quebrar'', ``decompor'', como é possível inferir e.g.~nos termos \emph{hidrólise} --- reação química de ``quebra'' da molécula da água em hidrogênio e oxigênio ---, \emph{pirólise} --- decomposição de moléculas através de reações químicas a altas temperaturas (``fogo'') ---, \emph{hemólise} --- decomposição de hemácias (células sanguíneas) ---, etc.} mais detida. Tal processo de abstração, ao abstrair ``nesse primeiro momento da análise o complexo de interações da categoria'' para focar a atenção ``nos atributos mais gerais de tal elemento'' específico, ``leva do concreto imediato até a síntese abstrata e ainda `rarefeita' de elementos da totalidade'', o que permite uma compreensão não só das características deste complexo, como também a elucidação de suas inter-relações com os demais complexos da totalidade \parencite[p. 138]{Fortes2013}. Dessa forma, esse processo de abstração busca separar, idealmente, ``de maneira provisória toda `influência secundária' `estranha' a este momento específico'' que está sendo analisado \parencite[p. 60]{Fortes2016a}.

A compreensão de um complexo específico enquanto momento de uma totalidade requer que eventualmente ``retorne-se'' à visão deste todo. Neste ``caminho de volta'', deve-se considerar, agora explicitamente, determinações até agora abstraídas, em particular as interações com outros complexos dentro desta totalidade em que nosso objeto de estudo preliminar está inserido. Dessa forma, esta totalidade aparece agora ``não como a representação caótica de um todo, mas como uma rica totalidade de muitas determinações e relações'' \parencite[p. 54]{Marx2011}.

Destaque-se, por completude, que este procedimento ideal, de compreensão mais aprofundada da interação de complexos parciais com os demais complexos dentro de uma totalidade, não é a mesma coisa que essas interações em si. É por isso que Marx distingue a realidade concreta, da qual a análise ideal parte, do ``concreto no pensamento'', pois este ``retorno'' da análise abstrata rumo à totalidade que fora abstraída ``é somente o modo de pensamento de apropriar-se do concreto, \emph{de reproduzi-lo como um concreto mental}. Mas \emph{de forma alguma é o processo de gênese do próprio concreto}'' \parencite[p. 54-5, grifo meu]{Marx2011}.\footnote{É neste sentido também que Marx, em sua famosa passagem do prefácio da segunda edição d'\emph{O Capital}, em que afirma que, ao contrário do ``método de Hegel'', em ``seu método'' ``o ideal não é mais do que \emph{o material, transposto e traduzido na cabeça do homem}'' \parencite[p. 90, grifo meu]{Marx2017}.}

Neste presente trabalho, o objeto analisado pressupõe o modo de produção capitalista enquanto totalidade, e sua análise requer que se tenha em vista os complexos constitutivos desta totalidade. Restrinjo-me neste presente trabalho, porém, à análise da assim chamada ``esfera econômica'', abstraindo de suas inter-relações com outros complexos da totalidade social, dentre os quais vale destacar a esfera do Estado. (Afinal, há mais coisas entre o céu e a terra do que permite o escasso tempo de um programa de mestrado.) Tal escolha, porém, não é arbitrária.

O objetivo da presente discussão é mostrar que a própria lógica interna do capital conduz a um aumento do consumo energético, mas \emph{não} --- como os economistas neoclássicos tendem a afirmar --- porque o consumo ``\emph{residencial}'' aumenta.\footnote{O consumo energético residencial contempla todo uso de energia dentro dos domicílios (\emph{households} em inglês). Em termos macroeconômicos, diz respeito a uma parcela do ``consumo agregado das famílias''.} Pelo contrário, o aumento deste consumo é uma \emph{consequência} ulterior, justamente porque este, assim como todo consumo energético, é consumo \emph{de mercadorias} e, assim como nosso pão de cada dia não vem a nós puramente pela benevolência do padeiro, nossas fontes de energia útil, na forma de carvão, lenha, eletricidade etc., elas sempre têm uma etiqueta de preços. Em suma: os ``serviços energéticos'', usufruídos residencialmente ou não, não surgem espontaneamente no mercado: como todos os ``bens e serviços'', têm de ser \emph{produzidos}. É justamente porque o ``primeiro ato histórico'' consiste na reprodução do ser humano enquanto gênero --- que pressupõe, claro, sua reprodução enquanto indivíduo singular --- que a esfera da \emph{produção} possui prioridade ontológica na totalidade da vida humana.\footnote{``\ldots{} devemos começar por constatar o primeiro pressuposto de toda a existência humana e também, portanto, de toda a história, a saber, o pressuposto de que os homens têm de estar em condições de viver para poder `fazer história'. Mas, para viver, precisa-se, antes de tudo, de comida, bebida, moradia, vestimenta e algumas coisas mais. O primeiro ato histórico é, pois, a produção dos meios para a satisfação dessas necessidades, a produção da própria vida material, e este é, sem dúvida, um ato histórico, uma condição fundamental de toda a história, que ainda hoje, assim como há milênios, tem de ser cumprida diariamente, a cada hora, simplesmente para manter os homens vivos.'' \parencite[p. 32-3]{Marx2007}}

Já se pode entrever, portanto, que o consumo ``produtivo'' de energia tem um papel, no mínimo, não-desprezível no que diz respeito ao objeto do presente trabalho, pois toda produção demanda energia,\footnote{Aqui tomado no sentido da energia enquanto um bem/serviço a se usufruir, não no sentido físico e trivial de que ``tudo é energia''.} em maior ou menor grau. Isso se dá porque a esfera do consumo --- ou, em geral, a esfera da circulação --- depende \emph{ontologicamente} da esfera da produção: o que se pode consumir está delimitado pelo que é produzido, e é somente através do aumento da produção que um aumento do consumo se faz possível. Da forma forma, embora não seja (totalmente) verdade que a oferta crie sua própria demanda, ao menos é certo que a diminuição de uma oferta pré-existente engendrará uma frustração de parte da demanda existente.

Ora, o caráter imanente do capital, de valorização do valor, demanda que a produção de mercadorias aumente. Através da constatação de que o emprego da energia torna-se quintessencial à dinâmica do capital, concluir-se-á neste presente trabalho que o postulado de Khazzoom-Brookes é uma lei de tendência própria do modo de produção capitalista --- embora não exatamente no sentido proposto pelos economistas que abordam este fenômeno. Dito de outra forma, é um dos corolários necessários do modo de produção capitalista, mesmo que se trate de uma \emph{lei de tendência} que não necessariamente se efetive a todos os momentos.

\hypertarget{leis-de-tenduxeancia}{%
\section{Leis de tendência}\label{leis-de-tenduxeancia}}

Como excelentemente sumarizado por Prado, o objetivo da crítica ontológica do modo de produção capitalista (e de seus fenômenos socioeconômicos) é, como na tradição do realismo crítico,

\begin{quote}
descobrir os poderes, os mecanismos, as tendências inscritas nas estruturas {[}sociais etc.{]}, os quais governam o curso dos eventos. Explanar {[}no sentido de `crítica explanatória' (NFV){]} nessa visão, pois, vem a ser desvendar as leis tendenciais que se combinam, ou seja, se somam ou se opõem, na regulação dos fenômenos. \parencite[p. 40]{Prado2009}
\end{quote}

Marx afirma, no prefácio da primeira edição d'\emph{O Capital}, que o conteúdo de sua obra ``não se trata do grau maior ou menor de desenvolvimento dos antagonismos sociais decorrentes de leis naturais da produção capitalista. Trata-se \emph{dessas próprias leis}, dessas \emph{tendências que atuam e se impõem com férrea necessidade}'' \parencite[p. 78, grifo meu]{Marx2017}. A noção de \emph{leis de tendência} é um conceito crucial ao presente trabalho, cuja compreensão, porém, é beneficiada pela apreensão do conceito de ``retrodução''.

A retrodução é um procedimento lógico através do qual fenômenos empíricos são vistos enquanto ``resultado{[}s{]} sintético{[}s{]} da operação articulada de múltiplas causas que povoam o mundo como processos duradouros ({[}i.e.{]} tendências {[}\ldots{]}) capazes de produzir efeitos empiricamente constatáveis'' \parencite[p. 83]{Medeiros2021a}. O exemplo fundamental de tendência, para nossos propósitos, é o valor das mercadorias, que reconhecemos ``\emph{por detrás do movimento dos preços}'' (ibid., grifo meu). O procedimento retrodutivo busca obter a essência por trás deste movimento, que faz com que esta seja sua forma necessária de manifestação empírica.

Não se trata, porém, de uma relação entre fenômenos empíricos e leis causais empíricas. Trata-se de um procedimento ontologicamente motivado, de aferição das tendências \emph{não-empíricas} que jazem por detrás dos fenômenos empíricos, as quais ditam suas especificidades e aparência necessária. Como excelentemente resumido por Medeiros e Bonente:

\begin{quote}
Com efeito, a concretude da realidade efetivada tem por pressuposto a abstrata possibilidade de efetivar-se; a concretude do singular tem por pressuposto a abstrata universalidade; a concretude de um momento determinado da existência tem por pressuposto a malha abstrata de relações que determinam sua posição e importância na totalidade de que é parte. Ser concreto ou abstrato não é, em suma, ser mais ou menos complexo ou real, mas, para repetir, estar mais ou menos distante da realidade efetivada como possibilidade realizada. \parencite[p. 86-7]{Medeiros2021a}
\end{quote}

Elucida-se, assim, o procedimento da abstração acima abordado: busca-se, através da ``força da abstração'', inferir as tendências e legalidades imanentes da totalidade concreta analisada, com o que esta totalidade aparece agora ``não como a representação caótica de um todo, mas como uma rica totalidade de muitas determinações e relações'' \parencite[p. 54]{Marx2011}, em que os fenômenos socioeconômicos, antes vistos como acontecimentos singulares e desconexos, aparecem agora como consequências de interações de certos complexos e legalidades, isto é, aparecem como possuidores de ``determinação e especificidade concreta do tipo de lei que se afirma na interseção das interações'' de seus complexos originadores \parencite[p. 338]{Lukacs2012}.

Embora seja possível argumentar que o objetivo da Ciência como um todo é buscar os genuínos mecanismos causais dos fenômenos, o desafio das ciências humanas, e da Economia em particular, é, invés de buscar descrever as leis naturais independentes do fazer humano (como nas ciências exatas), ter como objeto um complexo de complexos cujas leis de tendência são \emph{engendradas pela própria ação humana}.

Este desafio, porém, aparece como possibilidade, pois a crítica de tais estruturas sociais, ao desnaturalizá-las, elucida sua temporalidade, e põe em questão a possibilidade de sua superação, em última instância, política. Como o próprio Marx o coloca, em sua configuração racional, o caráter dialético da crítica da economia política ``constitui um escândalo e um horror para a burguesia e seus porta-vozes doutrinários, uma vez que, na intelecção positiva do existente, inclui, ao mesmo tempo, a intelecção de sua negação, de seu necessário perecimento'', assim como ``não se deixa intimidar por nada e é, por essência, crítica e revolucionária'' \parencite[p. 91]{Marx2017}.

Pode parecer que a discussão no presente trabalho trilha um caminho demasiado tortuoso em direção à verdade. Mas, assim como em Álvaro de Campos ``as cartas de amor, se há amor, têm de ser ridículas; se não fossem ridículas, não seriam cartas de amor'', os caminhos em direção à verdade, se buscam de fato a verdade, têm de ser tortuosos, pois, se não fossem tortuosos, não seriam caminhos em direção à compreensão verdadeira do ``garrido tumulto dos acontecimentos''.\footnote{\textcite{Hegel1995}, p.~34.}

%% \hl{Poder-se-ia argumentar que este é o objetivo da Ciência como um todo: buscar os genuínos mecanismos causais dos fenômenos, num reiterado esforço de retrodução.\footnote{``A retrodução é um tipo de raciocínio lógico-filosófico (passível de ser transposto para campo científico) que embasa o pensamento crítico-realista. Ela formula um questionamento que visa reconstituir a trajetória de formação de uma entidade a partir de seu estágio atual, em direção ao passado. Por exemplo: \emph{Como teria que ser X para que se tornem objetos de conhecimento para nós? Que propriedades teriam o objeto X para que ele se tornasse objeto de conhecimento para nós (fosse possível ser conhecido)?}'' \parencite[p. 215, n. 4]{Pimentel2017}}}
